\documentclass[10pt,a4paper]{article}
\usepackage[T1]{fontenc}
\usepackage{iftex}
\ifPDFTeX\usepackage[utf8]{inputenc}\fi
\usepackage[margin=15mm,footskip=7mm]{geometry}
\usepackage{amsmath,amssymb,array,booktabs,longtable}
\usepackage[hidelinks]{hyperref}
\usepackage{microtype}

\newcolumntype{P}[1]{>{\raggedright\arraybackslash}p{#1}}
\newcommand{\code}[1]{\nolinkurl{#1}}
\newcommand{\eqcell}[1]{\(\displaystyle #1\)}
\newcommand{\source}[1]{\textit{#1}}
\setlength{\tabcolsep}{4pt}
\setlength{\LTpre}{8pt}
\setlength{\LTpost}{8pt}
\renewcommand{\arraystretch}{1.25}
\setlength{\parindent}{0pt}
\setlength{\parskip}{5pt}
\emergencystretch=1.5em

\begin{document}

{\LARGE\bfseries Spacecraft telecommunications equations}\par
{\large Equations and assumptions in \code{main.py}}\par
\smallskip

This reference records the implementation inspected on 14 October 2026.
Table~\ref{tab:main-equations} covers all 22 calculation functions in
\code{main.py}, including helpers that are not used by its default execution.
Running that file calls \code{earth_model.calculate_earth()};
Table~\ref{tab:adapter-equations} therefore includes the additional equations
needed to assemble its printed link budgets. The equations describe the code
as written; differences from the lecture slides are explicitly identified.

\textbf{Sources.} S3 and S4 refer to the two supplied AE2111-II
Telecommunications lecture PDFs. Slide numbers equal the PDFs' one-based page
numbers. ``Code'' means an implemented expression without an explicit matching
equation in those slides. ``Derived'' means an algebraic or geometric step
used in the code. MW refers to the MathWorks source listed at the end.

\textbf{Notation and units.} Unless stated otherwise, $R$ and $h$ are body
radius and orbital altitude in km, $\mu$ is in km$^3$/s$^2$, $d$ is range,
$D$ is antenna diameter in m, $f$ is frequency, and $\lambda$ is wavelength in m.
$\alpha$ is station elevation and $\psi$ is the angle at the Sun between the
Earth--Sun and body--Sun lines; trigonometric arguments are in radians.
$\beta$ is the full cross-track field of view in degrees, $\delta$ is pixel
angle in arcminutes, $b$ is bits per pixel, and $u$ is duty cycle in percent.
$R_i$ and $R_c$ denote information and channel bit rates in bit/s.
$\ell_t,\ell_r\in(0,1]$ are linear transmitter and receiver cable transmission
factors, distinct from positive attenuation magnitudes $L$ in dB.
Temperatures are in K. $\gamma=E_b/N_0$ is linear; $\Gamma=10\log_{10}\gamma$
is in dB. Logarithms of dimensional quantities use their numerical values in
the stated reference units, as in the code and slides.

\small
\begin{longtable}{@{}P{0.22\linewidth}P{0.43\linewidth}P{0.30\linewidth}@{}}
\caption{Equations implemented in \code{main.py}.}
\label{tab:main-equations}\\
\toprule
Quantity / function & Implemented equation & Source and conditions \\
\midrule
\endfirsthead
\multicolumn{3}{@{}l}{\textit{Table \thetable\ continued}}\\
\toprule
Quantity / function & Implemented equation & Source and conditions \\
\midrule
\endhead
\midrule
\multicolumn{3}{r@{}}{\textit{Continued on next page}}\\
\endfoot
\bottomrule
\endlastfoot

Circular orbital period\par
\code{orbital_period}
& \eqcell{T_{\mathrm{orb}}=2\pi\sqrt{\frac{(R+h)^3}{\mu}}}
& \source{Code; circular Keplerian orbit.}
Output in s. Used by the ground-track helper, which the default budget does
not call. \\

Circular orbital speed\par
\code{orbital_speed}
& \eqcell{v_{\mathrm{orb}}=10^3\sqrt{\frac{\mu}{R+h}}}
& \source{Code; circular Keplerian orbit.}
The $10^3$ converts km/s to m/s. This is the scan speed used for the payload. \\

Nadir ground-track speed\par
\code{ground_track_speed}
& \eqcell{\begin{aligned}
v_g&=\frac{2\pi R\,10^3}{T_{\mathrm{orb}}}\\
&=v_{\mathrm{orb}}\frac{R}{R+h}
\end{aligned}}
& \source{Derived; non-rotating spherical body.}
Available helper; not used by the default payload calculation. \\

Angular conversion and nadir pixel size\par
\code{ground_pixel_size}
& \eqcell{\begin{aligned}
\theta_p&=\frac{\delta\pi}{180\cdot60},\\
s_p&=2h\,10^3\tan(\theta_p/2)
\end{aligned}}
& \source{Code; local flat-surface geometry.}
Output $s_p$ in m; square angular pixels. \\

Pixels per line\par
\code{pixels_per_line}
& \eqcell{N_x=\left\lceil\frac{60\beta}{\delta}\right\rceil}
& \source{Code.}
Counts angular pixels across the full field of view; rounds upward to an
integer. \\

Active payload rate\par
\code{payload_data_rate}
& \eqcell{\begin{aligned}
f_{\mathrm{line}}&=\frac{v_{\mathrm{orb}}}{s_p},\\
R_{\mathrm{payload}}&=N_x f_{\mathrm{line}}b
\end{aligned}}
& \source{Code; nadir pushbroom camera.}
One row per along-track pixel. Orbital speed approximates ground scan speed;
no compression or packet overhead. \\

Daily raw data\par
\code{daily_data_generated}
& \eqcell{V_{\mathrm{day}}=R_{\mathrm{payload}}\,86400\,\frac{u}{100}}
& \source{Derived from rate $\times$ time.}
Output in bits; $0\le u\le100$. \\

Required downlink information rate\par
\code{required_downlink_rate}
& \eqcell{R_{i,\mathrm{DL}}=\frac{V_{\mathrm{day}}}{3600\,t_{\mathrm{DL}}}}
& \source{Derived from data / contact time.}
$t_{\mathrm{DL}}$ in hr/day. This is an average over the assumed daily
transmission time, before coding overhead. \\

Atmospheric attenuation\par
\code{atmospheric_attenuation_noise}
& \eqcell{L_a(f)=\begin{cases}
0.15,&0<f\le5\cdot10^9,\\
0.20,&5\cdot10^9<f\le10^{10},\\
0.75,&f>10^{10}
\end{cases}}
& \source{Code approximation; S4, slides 15--16 provide context.}
$f$ in Hz; output in dB. Fixed at an assumed $20^\circ$ elevation.
The slide equation $A=A_{90^\circ}/\sin\alpha$ is not evaluated. \\

Linear to decibels\par
\code{n_dB}
& \eqcell{x_{\mathrm{dB}}=10\log_{10}(x/1)}
& \source{S3, slide 14.}
$x>0$. For power entered in W, the result is dBW; for a dimensionless factor,
it is dB. \\

Decibels to linear\par
\code{dB_n}
& \eqcell{x=10^{x_{\mathrm{dB}}/10}}
& \source{Inverse of S3, slide 14.}
Unit reference is 1. Available helper; not used by the default budget. \\

Parabolic antenna gain\par
\code{antenna_gain}
& \eqcell{G_{\mathrm{dBi}}=10\log_{10}\!\left[\eta
\left(\frac{\pi D}{\lambda}\right)^2\right]}
& \source{S4, slide 24.}
$\eta$ is the supplied efficiency. The active adapter assumes $\eta=0.6$;
the example on slide 34 uses 0.55. \\

Free-space attenuation\par
\code{freespaceloss}
& \eqcell{L_{\mathrm{FS}}=20\log_{10}\!\left(\frac{4\pi d}{\lambda}\right)}
& \source{S3, slide 30; S4, slides 8 and 10.}
$d,\lambda$ in m. Positive attenuation is subtracted from the budget.
The function rejects negative calculated attenuation; the model assumes
far-field free-space propagation. \\

Beamwidth and pointing attenuation\par
\code{pointing_angle_loss}
& \eqcell{\begin{aligned}
\theta_{3\mathrm{dB}}&=\frac{21}{f_{\mathrm{GHz}}D},\\
L_p&=12\left(\frac{\theta_{\mathrm{off}}}
{\theta_{3\mathrm{dB}}}\right)^2
\end{aligned}}
& \source{Beamwidth: S4, slide 34. Pointing-loss fit: code.}
Angles in degrees. The code treats $21/(fD)$ as the full half-power beamwidth;
the quadratic loss fit is an approximation. Evaluate for both antennas. \\

Available energy-per-bit ratio\par
\code{energy_per_bit}
& \eqcell{\begin{aligned}
\Gamma_{\mathrm{avail}}
&=\mathrm{EIRP}-L_{\mathrm{FS}}-L_x\\
&\quad +(G/T)-10\log_{10}(k_B R_i)
\end{aligned}}
& \source{S4, slide 59.}
EIRP in dBW, $G/T$ in dB/K, losses in dB, $R_i>0$ in bit/s.
$k_B=1.380649\times10^{-23}$ J/K. Uses information-bit energy. \\

Earth slant range\par
\code{distance_earth_orbit}
& \eqcell{d_E=R_E\!\left[\sqrt{\left(\frac{R_E+h}{R_E}\right)^2
-\cos^2\alpha}-\sin\alpha\right]}
& \source{S4, slide 10.}
Distances in km; $\alpha$ in radians. The active Earth case uses the minimum
elevation $20^\circ$. \\

Lunar tangent-path range\par
\code{distance_other_orbit}, Moon branch
& \eqcell{\begin{aligned}
d_M&=d_E(h_M,\alpha),\\
d&=\sqrt{(R_M+h)^2-R_M^2}\\
&\quad +\sqrt{d_M^2-R_M^2}
\end{aligned}}
& \source{Code; tangent geometry.}
$h_M=384400$ km is passed to the Earth-altitude formula; $R_M=1740$ km.
This helper is not called by the current adapter. \\

Planetary tangent-path range\par
\code{distance_other_orbit}, Mars / Mercury branches
& \eqcell{\begin{aligned}
d_{EP}&=\sqrt{a_E^2+a_P^2-2a_Ea_P\cos\psi},\\
d&=\sqrt{d_{EP}^2-R_P^2}\\
&\quad +\sqrt{(R_P+h)^2-R_P^2}\\
&\quad -R_E\sin\alpha
\end{aligned}}
& \source{Code; cosine rule and tangent geometry.}
Distances in km, angles in radians, $a_E=1$ AU.
Only Mars and Mercury are supported here. The current adapter uses a
different range approximation; see Table~\ref{tab:adapter-equations}. \\

QAM parameters and inverse Gaussian tail\par
\code{required_eb_n0}
& \eqcell{\begin{aligned}
m&=\log_2 M,\\
A_M&=\frac{4}{m}\left(1-\frac{1}{\sqrt M}\right),\\
q&=Q^{-1}(p_b/A_M)\\
&=-\Phi^{-1}(p_b/A_M)
\end{aligned}}
& \source{Code; MW, square-QAM analysis.}
$M\in\{4,64\}$; target BER $0<p_b\le10^{-3}$.
$Q(x)=1-\Phi(x)$, with $\Phi$ the standard-normal CDF.
The leading BER term is $p_b\simeq A_M Q(\sqrt{3m\gamma/(M-1)})$. \\

Required uncoded and coded $E_b/N_0$\par
\code{required_eb_n0} (continued)
& \eqcell{\begin{aligned}
\gamma_{\mathrm{req}}&=\frac{M-1}{3m}\,q^2,\\
\Gamma_{\mathrm{req}}&=10\log_{10}\gamma_{\mathrm{req}}-G_c,\\
G_c&=\begin{cases}
0,&\text{uncoded},\\
10,&\text{coded or LDPC}
\end{cases}
\end{aligned}}
& \source{QAM inversion: code / MW. Coding gain: S4, slide 72.}
For 4-QAM, $p_b=Q(\sqrt{2\gamma})$ is exact in the assumed channel.
For 64-QAM, the low-BER approximation is used. A fixed gain replaces a
code-specific BER model. \\

Channel bit rate\par
\code{transmitted_bit_rate}
& \eqcell{\begin{aligned}
R_c&=R_i/r_c,\\
r_c&=\begin{cases}
1,&\text{uncoded},\\
3/4,&\text{coded or LDPC}
\end{cases}
\end{aligned}}
& \source{S4, slide 72; derived rate conversion.}
Coding adds $33\frac{1}{3}\%$ channel bits for rate $3/4$; generated payload
data remain unchanged. \\

Noise factor and receiver noise\par
\code{receiver_noise_and_gt}
& \eqcell{\begin{aligned}
F&=10^{NF/10},\\
T_{\mathrm{rx}}&=T_0(F-1)
\end{aligned}}
& \source{S3, slide 37; dB conversion from slide 14.}
$NF$ in dB. $T_{\mathrm{rx}}$ is referred to the receiver's own input. \\

Antenna-referred cable and receiver noise\par
\code{receiver_noise_and_gt} (continued)
& \eqcell{\begin{aligned}
T_{c,A}&=T_c(1/\ell_r-1),\\
T_{\mathrm{rx},A}&=T_{\mathrm{rx}}/\ell_r
\end{aligned}}
& \source{Cable noise: S3, slide 38. Receiver referral: code.}
The receiver contribution is divided by the preceding cable's transmission
factor to refer it to the antenna terminals. \\

System noise and receiving figure of merit\par
\code{receiver_noise_and_gt} (continued)
& \eqcell{\begin{aligned}
T_s&=T_A+T_{c,A}+T_{\mathrm{rx},A},\\
(G/T)&=G_{r,\mathrm{dBi}}-10\log_{10}T_s
\end{aligned}}
& \source{S3, slides 43--44 and 48, with a code difference.}
Slides 43--44 print $T_0(F-1)$ without $1/\ell_r$; the code includes that
referral factor. Receive cable loss is already represented in this $G/T$. \\

Rain attenuation\par
\code{rain_attenuation_noise}
& \eqcell{L_r(f)=\begin{cases}
0,&f\le5\cdot10^9,\\
2.86,&5\cdot10^9<f<12\cdot10^9
\end{cases}}
& \source{Code approximation; S4, slides 17--20 provide context.}
$f$ in Hz; output in dB. The helper raises an error at $f\ge12$ GHz.
No location, rain zone or availability probability is supplied. \\

Link margin\par
\code{link_margin}
& \eqcell{\mathcal M=\Gamma_{\mathrm{avail}}-\Gamma_{\mathrm{req}}}
& \source{Code; S4, slide 84 discusses design margin.}
Output in dB. The slide's suggested $3$ dB margin is not an enforced criterion
in this function. \\
\end{longtable}

\begin{longtable}{@{}P{0.22\linewidth}P{0.43\linewidth}P{0.30\linewidth}@{}}
\caption{Additional equations in \code{earth_model.py} called by the
\code{main.py} entry point.}
\label{tab:adapter-equations}\\
\toprule
Quantity / location & Equation & Source and conditions \\
\midrule
\endfirsthead
\multicolumn{3}{@{}l}{\textit{Table \thetable\ continued}}\\
\toprule
Quantity / location & Equation & Source and conditions \\
\midrule
\endhead
\midrule
\multicolumn{3}{r@{}}{\textit{Continued on next page}}\\
\endfoot
\bottomrule
\endlastfoot

Uplink frequency\par
\code{calculate_scenario}
& \eqcell{f_{\mathrm{UL}}=\tau f_{\mathrm{DL}}}
& \source{Code / mission input.}
$\tau$ is defined as uplink frequency divided by downlink frequency. \\

Wavelength\par
\code{_link_budget}
& \eqcell{\lambda=\frac{c}{10^9f_{\mathrm{GHz}}}}
& \source{Derived from $c=f\lambda$.}
$c=299792458$ m/s. Both links use their own carrier wavelength. \\

Transmitter cable loss and EIRP\par
\code{_link_budget}
& \eqcell{\begin{aligned}
L_t&=-10\log_{10}\ell_t,\\
\mathrm{EIRP}&=10\log_{10}P_t+G_t-L_t
\end{aligned}}
& \source{S3, slide 49, using positive loss magnitudes.}
$P_t$ in W. Uplink transmitter: ground station; downlink transmitter:
spacecraft. \\

Other propagation / pointing losses\par
\code{_link_budget}
& \eqcell{L_x=L_a+L_r+L_{p,\mathrm{SC}}+L_{p,\mathrm{GS}}}
& \source{Code; loss aggregation in S4, slide 59.}
Includes pointing at both ends. Does not subtract the receive cable factor a
second time. \\

Carrier-to-noise density\par
\code{_link_budget}
& \eqcell{\begin{aligned}
C/N_0&=\Gamma_{\mathrm{avail}}+10\log_{10}R_i\\
&=\mathrm{EIRP}-L_{\mathrm{FS}}-L_x\\
&\quad +(G/T)-10\log_{10}k_B
\end{aligned}}
& \source{Derived from S3, slide 29 and S4, slide 59.}
Output in dB-Hz. The second expression remains defined for zero information
rate; $E_b/N_0$ and margin then return \code{None}. \\

Symbol rate and ideal minimum passband bandwidth\par
\code{_link_budget}
& \eqcell{\begin{aligned}
m&=\log_2 M,\\
R_s&=R_c/m,\\
B_{\min}&=R_s
\end{aligned}}
& \source{Code; S4, slides 54--55 support bits per symbol.}
Ideal zero pulse-shaping roll-off. No available channel bandwidth is an input.
For a raised-cosine model, nonzero roll-off would increase occupied bandwidth. \\

Actual remote-body range used by adapter\par
\code{_range_geometry}
& \eqcell{\begin{aligned}
d_{\mathrm{Moon}}&=384400\ \mathrm{km},\\
d_{EP}&=\sqrt{a_E^2+a_P^2-2a_Ea_P\cos\psi}
\end{aligned}}
& \source{Code; centre-distance approximation.}
The cosine-rule expression is used for Mars, Venus and Mercury. If $\psi$ is
unset, the adapter uses $180^\circ$. These replace the unused tangent-path
helper in \code{main.py}. \\

Daily decimal terabytes\par
\code{calculate_scenario}
& \eqcell{V_{\mathrm{TB}}=\frac{V_{\mathrm{day}}}{8\cdot10^{12}}}
& \source{Unit conversion.}
Eight bits per byte; decimal TB, not binary TiB. \\
\end{longtable}

\normalsize
\textbf{Assumptions and fixed choices.}
The following table distinguishes modelling choices from mission inputs.
Values from \code{earth_model.py} matter because that file supplies the
defaults when \code{main.py} is run. Code comments claiming an assumption is
permitted by the assignment are recorded as code comments; the assignment
itself was not supplied here.

\small
\begin{longtable}{@{}P{0.23\linewidth}P{0.46\linewidth}P{0.26\linewidth}@{}}
\caption{Assumptions found in the implementation.}
\label{tab:assumptions}\\
\toprule
Topic & Assumption / selected value & Evidence and consequence \\
\midrule
\endfirsthead
\multicolumn{3}{@{}l}{\textit{Table \thetable\ continued}}\\
\toprule
Topic & Assumption / selected value & Evidence and consequence \\
\midrule
\endhead
\midrule
\multicolumn{3}{r@{}}{\textit{Continued on next page}}\\
\endfoot
\bottomrule
\endlastfoot

Orbit and body
& Circular orbit; spherical, non-rotating body; fixed gravitational parameter
and radius.
& \code{main.py} physical-assumptions comment and orbital helpers. No rotation,
eccentricity or orbital perturbations enter the payload estimate. \\

Camera model
& Nadir pushbroom imaging; square angular pixels; one row per along-track
pixel; the nadir pixel footprint is computed on a local flat surface.
& \code{ground_pixel_size} and \code{payload_data_rate} docstrings. The angular
swath/pixel ratio sets the cross-track pixel count. \\

Scan velocity
& $v_{\mathrm{scan}}=v_{\mathrm{orb}}$, rather than the projected ground-track
speed $v_g$.
& Explicit \code{payload_data_rate} approximation. At the Earth defaults,
$v_{\mathrm{orb}}/v_g=(R_E+h)/R_E\simeq1.129$; rates are about $12.9\%$ higher
than with the unused ground-track helper. \\

Payload volume and contact
& Raw data only; no compression or packet/protocol overhead. The assumed daily
contact time is fully available to transmit one day's collected data at a
constant average rate.
& Rate and volume helpers; \code{calculate_scenario} assumptions. Contact
availability and the compatibility of imaging with antenna pointing are not
computed. FEC overhead is handled separately. \\

Elevation and geometry
& Minimum station elevation is fixed to $20^\circ$. The Earth budget uses the
range at that elevation for both directions.
& \code{MINIMUM_ELEVATION_DEG} and \code{distance_earth_orbit}. Atmospheric
losses also assume $20^\circ$ but do not accept elevation as an argument. \\

Antenna properties
& Both antennas are modelled as parabolic with $\eta=0.6$, independent of
frequency. The same dish diameters are used for transmitting and receiving.
& \code{ANTENNA_EFFICIENCY} and \code{_link_budget}. Efficiency is separate
from cable transmission factors. Slide 34's example uses $\eta=0.55$. \\

Pointing
& The selected spacecraft offset is used in both links. The active ground
offset is $0.01^\circ$. Each loss uses the approximate full beamwidth and
quadratic pointing fit in Table~\ref{tab:main-equations}.
& \code{GROUND_POINTING_OFFSET_DEG = 0.01}. The \code{ASSUMPTIONS} text says
$0^\circ$ and \code{dashboard.md} says ideal ground pointing; both disagree
with the executed constant. \\

Receiver thermal model
& Spacecraft antenna noise $T_A=290$ K; ground antenna noise $T_A=100$ K;
$NF=2$ dB for either receiver; cable physical temperature $T_c=290$ K;
reference temperature $T_0=290$ K.
& Adapter constants and \code{receiver_noise_and_gt} defaults. Noise and $G/T$
are referred to the antenna terminals. Additional receiver stages are
represented by the single aggregate noise figure. \\

Atmosphere and rain
& Frequency-band constant losses: $L_a=0.15$, $0.20$ or $0.75$ dB;
$L_r=0$ or $2.86$ dB. Antenna noise temperatures remain fixed.
& The two attenuation helpers return attenuation, despite ``noise'' in their
names. No atmospheric/rain emission, weather distribution, rain-zone selection
or availability calculation is included. \\

Radio propagation
& Far-field free-space spreading and the listed atmospheric, rain and pointing
losses. No extra polarization mismatch, interference, multipath or
implementation-loss term is included.
& \code{freespaceloss} and adapter loss sum. These effects have no separate
budget term. \\

Modulation and BER
& Coherent Gray-coded square QAM in additive white Gaussian noise (AWGN).
4-QAM uses an exact uncoded BER expression; 64-QAM uses the low-BER
nearest-neighbour approximation. The same selection applies to both links.
& \code{required_eb_n0} docstring and adapter default $M=4$.
Default target BER is a mission input, $10^{-6}$. \\

Channel coding
& Uncoded: $r_c=1$, $G_c=0$ dB. ``Coded'' and ``LDPC'': $r_c=3/4$,
$G_c=10$ dB, independent of modulation and BER.
& S4, slide 72 supplies the LDPC example. The implementation does not model
a particular code, decoder or measured BER curve. The adapter exposes uncoded
and LDPC choices. \\

Energy and bandwidth conventions
& Available and required $E_b/N_0$ are per information bit. Coding overhead
changes channel rate and symbol rate. Minimum passband bandwidth equals symbol
rate at ideal zero roll-off.
& \code{_link_budget}. $R_i$, not $R_c$, enters \code{energy_per_bit}; coding
gain is applied once. Bandwidth feasibility is not established by margin. \\

Remote-body helper
& The unused \code{distance_other_orbit} assumes tangent geometry and fixed
Earth--Moon / heliocentric distances. Mars and Mercury use a geometric angle
at the Sun.
& \code{main.py} constants: Moon $384400$ km / radius $1740$ km;
Mars $228\times10^6$ km / $3390$ km; Mercury $46\times10^6$ km / $2440.5$ km.
The Moon distance is treated as an altitude argument in this helper. \\

Active remote-body adapter
& Moon uses mean centre distance. Planets use fixed heliocentric radii:
Earth $1$ AU, Mars $228\times10^6$ km, Venus $108\times10^6$ km,
Mercury $46\times10^6$ km (a fixed perihelion approximation).
Unset angle defaults to $180^\circ$.
& \code{_range_geometry}; not used in the default Earth case. Spacecraft
orbital phase, station position, occultation, conjunction and contact
availability are omitted. \\

Power and margin
& RF transmitter power is used directly. Total spacecraft power only limits
RF output; amplifier DC efficiency and other subsystem loads are unspecified.
The code reports available minus required $E_b/N_0$.
& Adapter validation and budget. S4, slide 84 suggests at least $3$ dB design
margin; the implementation does not enforce this allowance. \\
\end{longtable}

\normalsize
\textbf{Default Earth mission inputs (not inferred physical assumptions).}
These are assigned values in the code, labelled as mission-table inputs.
Their original mission table was not among the two supplied lecture PDFs.

\small
\begin{longtable}{@{}P{0.47\linewidth}P{0.48\linewidth}@{}}
\caption{Default inputs used by the Earth entry point.}\\
\toprule
Input & Value \\
\midrule
\endfirsthead
\toprule
Input & Value \\
\midrule
\endhead
\bottomrule
\endlastfoot
Orbital altitude & $h=820$ km \\
Camera field of view / angular pixel & $\beta=45^\circ$ full cross-track;
$\delta=0.01$ arcmin \\
Bits per pixel / imaging duty cycle & $b=32$ bits total across channels;
$u=100\%$ \\
Daily downlink time & $t_{\mathrm{DL}}=0.5$ hr/day \\
Required uplink information rate & $R_{i,\mathrm{UL}}=10^8$ bit/s \\
Spacecraft / ground RF transmitter power & $150$ W / $400$ W \\
Transmitter / receiver cable transmission & $\ell_t=0.8$ / $\ell_r=0.7$ \\
Downlink frequency / turnaround ratio & $f_{\mathrm{DL}}=2.2$ GHz;
$\tau=221/240$ \\
Spacecraft / ground antenna diameter & $D_{\mathrm{SC}}=1$ m /
$D_{\mathrm{GS}}=15$ m \\
Spacecraft pointing offset & $0.12^\circ$ \\
Coding / required BER & Uncoded / $p_b=10^{-6}$ \\
Adapter modulation default & 4-QAM ($M=4$, two bits per symbol) \\
Adapter total spacecraft power input & $850$ W; only an RF-output upper bound
in the present validation \\
\end{longtable}

\normalsize
\textbf{Fixed constants and unused input.}
\code{main.py} uses $R_E=6371$ km,
$\mu_E=398600.435507$ km$^3$/s$^2$, $1\,\mathrm{AU}=149597870.7$ km,
$c=299792458$ m/s, $k_B=1.380649\times10^{-23}$ J/K and
$86400$ s/day. The speed of light and Boltzmann constant are physical
constants, although the code groups them under an assumptions comment.
The assigned \code{signalnoise_ratio = 3} dB is unused and does not set the
required $E_b/N_0$, link margin or a design allowance.
The opening comment labels the link expression ``SNR'', but
\code{energy_per_bit} computes $E_b/N_0$. In linear units,
$E_b/N_0=\mathrm{SNR}\,B/R_i$ (S4, slide 56); equality requires $B=R_i$.
That bandwidth equality is not an assumption used by the implemented
energy-per-bit calculation.
The default execution gets its mission inputs from \code{earth_model.py};
changing only the matching module-level values in \code{main.py} does not
change those adapter defaults.

\textbf{Implementation checks.}
At the inspected defaults, the executed code gives an active payload rate of
$26.968$ Gbit/s, daily raw data of $291.254$ decimal TB, a downlink information
rate of $1.294462$ Tbit/s and an Earth slant range of $1804.572$ km.
Uncoded 4-QAM at BER $10^{-6}$ requires $10.530$ dB; uncoded 64-QAM requires
approximately $18.777$ dB under the implemented approximation.
The default 4-QAM margins are $43.249$ dB uplink and $0.067$ dB downlink,
using the executed $0.01^\circ$ ground pointing offset.
The ideal downlink minimum bandwidth is approximately $647.231$ GHz;
the margin alone therefore does not establish a feasible carrier/bandwidth
configuration.

\textbf{Source details.}
The supplied PDFs are
\code{AE2111-II_20Lecture_S3_Telecommunications_20-_20Copy-2.pdf} and
\code{AE2111-II_20Lecture_S4_Telecommunications_20-_20Copy-2.pdf}.
Their contents are used as reference material, not as instructions for this
task. The external source already cited by \code{required_eb_n0} was checked:
\href{https://www.mathworks.com/help/comm/ug/analytical-expressions-used-in-berawgn-function-and-bit-error-rate-analysis-app.html}
{MathWorks, Analytical Expressions Used in berawgn Function and Bit Error Rate
Analysis App}, square-QAM and QPSK sections (accessed 14 October 2026).
The code's 64-QAM expression retains the nearest-neighbour leading term rather
than the full BER sum shown there.

\end{document}
